\documentclass[11pt,a4paper]{article}
\usepackage[margin=1in]{geometry}
\usepackage{amsmath,amssymb}
\usepackage{parskip}
\usepackage{hyperref}

\title{\textbf{Pre-Experiment Overview:\\ Atomic Spectra and Determination of Rydberg's Constant}}
\author{IBRAHIM H.I. ABUSHAWISH \\ Istanbul University, Department of Physics}
\date{}

\begin{document}
\maketitle

\section*{Purpose of the Experiment}
According to the lab manual, the primary purpose of this experiment is to explore the connection between atomic emission spectra and discrete atomic energy levels using hydrogen and mercury spectral lamps. By passing emitted light through a transmission diffraction grating, we measure the diffraction angles and wavelengths of visible spectral lines (notably the Balmer series of hydrogen: $H_\alpha$, $H_\beta$, and $H_\gamma$). From these spectral line measurements, we calculate the empirical Rydberg constant, $R_H$, and compare it with the theoretical quantum mechanical value.

\section*{Importance and Significance}
This experiment is a foundational pillar of modern quantum physics because discrete line spectra could never be explained by classical electrodynamics, which predicted that accelerating orbiting electrons should continuously radiate and collapse into the nucleus. Instead, Niels Bohr postulated that electrons inhabit stationary, non-radiating quantum orbits, and light is emitted only when an electron jumps from a higher level, $n_2$, to a lower level, $n_1$, with photon energy $\Delta E = h c / \lambda$.

In the laboratory setup, light emitted from excited gas atoms passes through a transmission grating with known grating constant $g$. The waves diffract and interfere constructively according to the grating equation $m \lambda = g \sin\theta$. By recording the distance between the grating and the viewing scale and locating the first-order and second-order diffraction positions, we isolate the exact wavelength of each visible line. Substituting these values into the Balmer-Rydberg relationship, $\frac{1}{\lambda} = R_H \left(\frac{1}{2^2} - \frac{1}{n^2}\right)$, allows us to extract Rydberg's constant with impressive precision. It is an extraordinary direct visualization of quantum jumps and orbital transitions right before your eyes.

\end{document}
